\documentclass[../../../include/open-logic-section]{subfiles}

\begin{document}

\olfileid[id]{sth}{card-arithmetic}{fix}
\olsection{Titik Tetap-$\aleph$}

Dalam \olref[spine][]{chap}, kita mengemukakan bahwa Penggantian memerlukan
pembenaran karena aksioma itu memaksa hierarki menjadi cukup tinggi. Setelah
mempelajari aritmetika kardinal, kita dapat memberikan suatu ilustrasi kecil
tentang ketinggian hierarki tersebut.

Jelas bahwa $0 < \aleph_0$, $1 < \aleph_1$, dan $2 < \aleph_2$\ldots{}
dan, sesungguhnya, selisih ukurannya semakin \emph{besar} pada setiap
langkah. Karena itu, kita mungkin tergoda untuk menduga bahwa
$\kappa<\aleph_\kappa$ bagi setiap ordinal~$\kappa$.

Namun, berdasarkan $\ZFC$, dugaan ini \emph{salah}. Bahkan, kita dapat
membuktikan bahwa terdapat \emph{titik-titik tetap-$\aleph$}, yakni kardinal
$\kappa$ sedemikian sehingga $\kappa=\aleph_\kappa$.

\begin{prop}\ollabel{alephfixed}
Terdapat suatu titik tetap-$\aleph$.
\end{prop}

\begin{proof}
Dengan menggunakan rekursi, definisikan:
\begin{align*}
	\kappa_0 &= 0\\
	\kappa_{n+1} &= \aleph_{\kappa_n}\\
	\kappa&= \bigcup_{n < \omega}\kappa_n
\end{align*}
Sekarang $\kappa$ merupakan kardinal berdasarkan
\olref[cardinals][classing]{unioncardinalscardinal}. Namun:
\[
	\kappa= \bigcup_{n < \omega} \kappa_{n+1} =
	\bigcup_{n < \omega}\aleph_{\kappa_n} =
	\bigcup_{\alpha < \kappa}\aleph_\alpha = \aleph_\kappa
\]
%	By construction, $\kappa$ is the least cardinal greater than each
%	$\kappa_n$. So $\aleph_\kappa$ is the least cardinal greater than
%	each $\aleph_{\kappa_n}$, and hence greater than each
%	$\kappa_{n+1}$. But equally $\kappa$ is the least cardinal greater
%	than each $\kappa_{n+1} = \aleph_{\kappa_n}$. So $\kappa=
%	\aleph_\kappa$.
\end{proof}

Boolos pernah menulis sebuah artikel tepat mengenai titik tetap-$\aleph$ yang
baru saja kita konstruksi. Setelah mencatat keberadaan $\kappa$ pada awal
artikelnya, ia berkata:
\begin{quote}
	[$\kappa$ merupakan] suatu bilangan yang \emph{cukup besar}, menurut
	ukuran orang-orang yang belum pernah mengenal teori himpunan; begitu
	besarnya, menurut saya, sehingga ia membuat kita mempertanyakan kebenaran
	setiap teori yang salah satu pernyataannya mengklaim bahwa terdapat
	sekurang-kurangnya $\kappa$ objek.
	\citep[hlm.~257]{Boolos2000}
\end{quote}
Pada akhirnya, ia menutup makalahnya dengan bertanya:
\begin{quote}
	[apakah] kita menduga bahwa, bagaimanapun keadaannya pada awal cerita,
	setelah mencapai sejauh ini roda-roda telah berputar di tempat dan kita
	tidak lagi mendengarkan deskripsi tentang sesuatu yang sungguh demikian
	adanya?
	\citep[hlm.~268]{Boolos2000}
\end{quote}
Jika kita memang telah melampaui ``segala sesuatu yang sungguh demikian
adanya'', maka kesalahan harus langsung ditimpakan kepada Penggantian. Sebab,
aksioma inilah yang memungkinkan bukti kita bekerja. Jika demikian, agaknya
Boolos perlu meninjau kembali klaim yang ia buat beberapa dasawarsa
sebelumnya bahwa Penggantian tidak mempunyai konsekuensi ``yang tidak
diinginkan'' (lihat \olref[replacement][extrinsic]{sec}).

Namun, apakah keberadaan $\kappa$ memang seburuk itu? Di sini mungkin
bermanfaat untuk mempertimbangkan \emph{paradoks Tristram Shandy} dari
Russell. Tristram Shandy mencatat kehidupannya dalam buku harian, tetapi ia
memerlukan satu tahun untuk mencatat satu hari. Setiap tahun berlalu,
Tristram semakin tertinggal: setelah satu tahun, ia baru mencatat satu hari
dan telah menjalani 364 hari yang belum tercatat; setelah dua tahun, ia baru
mencatat dua hari dan telah menjalani 728 hari yang belum tercatat; setelah
tiga tahun, ia baru mencatat tiga hari dan telah menjalani 1.092 hari yang
belum tercatat, \dots\footnote{Dengan mengabaikan tahun kabisat.} Meskipun
demikian, jika Tristram \emph{abadi}, ia akan berhasil mencatat setiap hari,
sebab ia akan mencatat hari ke-$n$ pada tahun ke-$n$ kehidupannya. Jadi,
``pada akhir waktu'', Tristram akan mempunyai buku harian yang lengkap.

Sekarang, mengapa hal ini begitu berbeda dari gagasan bahwa $\alpha$ lebih
kecil daripada $\aleph_\alpha$---dan bahkan semakin lama semakin jauh lebih
kecil---hingga mencapai $\kappa$, ketika kita akhirnya menyusul dan
$\kappa=\aleph_\kappa$?

Kesampingkan hal itu, dan andaikan kita menerima $\ZFC$. Mari kita akhiri
dengan sedikit kesenangan tambahan mengenai konstruksi titik tetap. Tiga
hasil berikut secara intuitif menetapkan bahwa terdapat suatu titik (yang
tidak trivial) ketika hierarki selebar tingginya:

\begin{prop}\ollabel{bethfixed}
Terdapat suatu titik tetap-$\beth$, yakni suatu $\kappa$ sedemikian sehingga
$\kappa=\beth_\kappa$.
\end{prop}

\begin{proof}
Seperti dalam \olref{alephfixed}, dengan menggunakan ``$\beth$'' sebagai
ganti ``$\aleph$''.
\end{proof}

\begin{prop}\ollabel{stagesize}
$\card{V_{\omega+\alpha}} = \beth_{\alpha}$. Jika $\omega \ordtimes
\omega \leq \alpha$, maka $\card{V_\alpha} = \beth_\alpha$.
\end{prop}

\begin{proof}
Klaim pertama berlaku berdasarkan suatu induksi transfinit sederhana. Klaim
kedua mengikuti karena, jika $\omega \ordtimes \omega \leq \alpha$, maka
$\omega + \alpha = \alpha$. Untuk menetapkannya, kita menggunakan fakta-fakta
tentang aritmetika ordinal dari \olref[ord-arithmetic][]{chap}. Pertama,
perhatikan bahwa $\omega \ordtimes \omega = \omega \ordtimes (1 \ordplus
\omega) = (\omega \ordtimes 1) \ordplus (\omega\ordtimes\omega) = \omega
\ordplus (\omega \ordtimes \omega)$. Sekarang, jika $\omega \ordtimes \omega
\leq \alpha$, yakni $\alpha = (\omega\ordtimes\omega) \ordplus \beta$ bagi
suatu $\beta$, maka $\omega \ordplus \alpha = \omega \ordplus ((\omega
\ordtimes \omega) \ordplus \beta) = (\omega \ordplus (\omega \ordtimes
\omega)) \ordplus \beta = (\omega \ordtimes \omega) \ordplus \beta =
\alpha$.
\end{proof}

\begin{cor}
Terdapat suatu $\kappa$ sedemikian sehingga $\card{V_\kappa} = \kappa$.
\end{cor}

\begin{proof}
Misalkan $\kappa$ suatu titik tetap-$\beth$ sebagaimana diberikan oleh
\olref{bethfixed}. Jelas bahwa $\omega \ordtimes \omega < \kappa$. Jadi,
$\card{V_\kappa} = \beth_\kappa= \kappa$ berdasarkan \olref{stagesize}.
\end{proof}

Terdapat sama banyak tahap di bawah $V_\kappa$ dengan !!{element}s dalam
$V_\kappa$. Jadi, secara intuitif, $V_\kappa$ selebar tingginya. Hal ini
sangat menyerupai Tristram Shandy: kita berpindah dari satu tahap ke tahap
berikutnya dengan mengambil \emph{himpunan-himpunan kuasa}, sehingga hierarki
menjadi \emph{jauh} lebih besar pada setiap langkah. Namun, ``pada akhirnya'',
yakni pada tahap $\kappa$, lebar hierarki menyusul tingginya.

Orang mungkin bertanya: \emph{Seberapa sering lebar hierarki menyamai
tingginya?} Jawabannya: \emph{Sesering adanya ordinal.} Namun, jawaban ini
memerlukan sedikit penjelasan.

Kita mendefinisikan suatu suku $\tau$ sebagai berikut. Untuk sebarang $A$,
tetapkan:
\begin{align*}
	\tau_0(A) & \defis \card{A}\\
	\tau_{n+1}(A) & \defis \beth_{\tau_n(A)}\\
	\tau(A) & \defis \bigcup_{n < \omega}\tau_n(A)
\intertext{Seperti dalam \olref{bethfixed}, $\tau(A)$ merupakan titik
tetap-$\beth$ bagi sebarang $A$, dan jelas $\card{A} < \tau(A)$. Sekarang
perhatikan definisi rekursif berikut:}
	W_0 &\defis 0\\
	W_{\alpha + 1} & \defis \tau(W_\alpha)\\
	W_\alpha & \defis \bigcup_{\beta < \alpha} W_\beta
	\text{, jika $\alpha$ merupakan ordinal limit}
\end{align*}
Konstruksi ini terdefinisi bagi semua ordinal. Jadi, secara intuitif, $W$
merupakan ``!!a{injection}'' dari ordinal-ordinal ke titik-titik tetap-$\beth$.
Dan, tepat seperti sebelumnya, $V_{W_\alpha}$ selebar tingginya bagi sebarang
$\alpha$.

\end{document}
